\documentclass[12pt]{article}
\usepackage[utf8]{inputenc}
\usepackage[T1]{fontenc}
\usepackage{fontspec}
\setmainfont{Times New Roman}
\newfontfamily\hackfont{Hack}
\usepackage{titlesec}
\usepackage{fancyhdr}
\usepackage{listings}
\usepackage{xcolor}
\usepackage{graphicx}
\usepackage{longtable}
\usepackage{caption}
\usepackage{hyperref}
\usepackage{setspace}
\usepackage{geometry}
\geometry{a4paper, margin=1in}
\hypersetup{colorlinks=true, linkcolor=blue, urlcolor=blue}
\renewcommand{\contentsname}{Table of Contents}

\pagestyle{fancy}
\fancyhf{} 							% Header Middle
\rhead{Computer Graphics Lab}
\lhead{Lab-02: Flag}
\rfoot{\thepage}

\begin{document}

\begin{titlepage}					% Cover Start
\begin{center}
    
\includegraphics[width=0.35\textwidth]{logo.png} \\
    \vspace{0.5cm}
    \textbf{\Huge \textcolor{cyan}{NOTRE DAME} \textcolor{cyan}{UNIVERSITY}} \\
    \vspace{0.5cm}
    \textbf{\Huge \textcolor{orange}{BANGLADESH}} \\
    
    \vspace{0.9cm}
    \textbf{\Huge \underline{Computer Graphics Lab Report-02}} \\
    
    \vspace{1.5em}
    \begin{flushleft}
    	\textbf{\Large Course Code: CSE-4204} \\
		\vspace{0.3cm}        
        \textbf{\Large Course Title: Computer Graphics Lab} \\
        \vspace{0.3cm}
        \textbf{\Large Lab Task Topic: Draw a Flag} \\        
    \end{flushleft}
    
    \vspace{1em}
    \begin{flushleft}
        \textbf{\Huge \textcolor{blue}{Submitted by:}} \\
        \vspace{0.3cm}
        \textbf{\Large Name: Istiak Alam} \\
        \vspace{0.3cm}
        \textbf{\Large ID: 0692230005101005} \\
		\vspace{0.3cm}        
        \textbf{\Large Batch: CSE-20} \\
        \vspace{0.5cm}
        \textbf{\Large Submission Date: }{\Large \textbf{\today}\par}
    \end{flushleft}
    \vfill
    \begin{flushleft}
        \textbf{\Huge \textcolor{blue}{Submitted to:}} \\
        \vspace{0.3cm}
        \textbf{\Large Humayara Binte Rashid} \\
        \vspace{0.3cm}
        \textbf{\Large Lecturer, Dept. of CSE} \\
        \vspace{0.3cm}
        \textbf{\Large Notre Dame University Bangladesh}
    \end{flushleft}
\end{center}
\end{titlepage}						% Cover End

\tableofcontents
\thispagestyle{empty}
\clearpage
\pagenumbering{arabic}

\section{Objective}
The objective of this lab is to understand and apply fundamental Computer Graphics concepts using OpenGL and GLUT. The experiment focuses on drawing a national flag by combining basic geometric primitives such as quadrilaterals, triangles, and circles within a 2D orthographic projection. Through this task, the use of coordinate systems, color models, and custom drawing functions is practiced to build a complete graphical object.


%-------------------------------------
\section{Tools and Environment}
\begin{itemize}
    \item Programming Language: C++
    \item Graphics Library: OpenGL with GLUT
    \item IDE: Code::Blocks
\end{itemize}


%-------------------------------------
\section{Graph Implementation}
Here is the Graph file \href{https://virtual-graph-paper.com/YmRiNDMzMjJlMTY2}{Link}
\begin{figure}[!h]
    \centering
    \includegraphics[width=0.75\textwidth]{graph.png}
    \caption{Graph View}
\end{figure}

%-------------------------------------
\section{Source Code}
The program starts by initializing the OpenGL environment, where the background color is set and a 2D orthographic projection is defined using the \texttt{glOrtho()} function. A custom function named \texttt{DrawCircle()} is implemented to draw the circular portion of the flag using the triangle fan technique and trigonometric calculations.

The \texttt{display()} function handles the rendering process. It draws the rectangular portion of the flag using quadrilateral primitives and places a filled circle at the center to represent the emblem. Additional geometric shapes are used to create the flag stand and decorative elements. Different colors are applied using the \texttt{glColor3f()} function to achieve a realistic appearance.

The \texttt{main()} function initializes the GLUT framework, creates the display window, registers the display callback function, and starts the main rendering loop to continuously display the flag. \\
\lstset{
    language=C++,
    basicstyle=\ttfamily\small,
    keywordstyle=\color{blue},
    commentstyle=\color{gray},
    stringstyle=\color{red},
    breaklines=true,
    frame=single
}

\begin{lstlisting}
/*
 * Lab Task-02
 * Creating a Flag
 * Written by Istiak Alam
 */

#include <GL/glut.h>  // Include GLUT header here
#include <stdlib.h>
#include <math.h>


//float p=-2.3;
void DrawCircle(float cx, float cy, float rx,float ry, int num_segments)
{
	glBegin(GL_TRIANGLE_FAN);
	for(int ii = 0; ii < num_segments; ii++)
	{
float theta = 2.0f * 3.1415926f * float(ii) / float(num_segments); //current angle

	float x = rx * cosf(theta);  //calculate the x component
	float y = ry * sinf(theta);  //calculate the y component
	
	glVertex2f(x + cx, y + cy);  //output vertex

	}
	glEnd();
}
void init(void)
{
	glClearColor (0.0, 0.0, 0.0, 0.0);
	 glOrtho(-20.0, 20.0, -20.0, 20.0, -1.0, 1.0);
}


void display(void)
{

    glClear(GL_COLOR_BUFFER_BIT);
	glPushMatrix();

    glColor3f(1.0f, 1.0f, 1.0f);  // Set color to white
    glBegin(GL_QUADS);  // Draw a square
        glVertex3d(2.0, 5.0, 0.0);
        glVertex3d(10.0, 5.0, 0.0);
        glVertex3d(10.0, 10.0, 0.0);
        glVertex3f(2.0, 10.0, 0.0);
    glEnd();

	glColor3f(0.7, 0.2, 0.3);
    DrawCircle(6,7.5f,1.5f,1.5f,100);
    glPopMatrix();

    glColor3f(0.3f, 0.0f, 0.0f);
    glBegin(GL_QUADS);  // Draw a square
        glVertex3d(1.0, -8.0, 0.0);
        glVertex3d(2.0, -8.0, 0.0);
        glVertex3d(2.0, 10.0, 0.0);
        glVertex3f(1.0, 10.0, 0.0);
    glEnd();

    glColor3f(0.9f, 1.0f, 0.3f);
        glBegin(GL_TRIANGLES);
        glVertex3d(1,10,0);
        glVertex3d(2,10,0);
        glVertex3d(1.5f,10.5f,0);
    glEnd();

	glFlush();
}

int main(int argc, char** argv)
{
    glutInit(&argc, argv);
	glutInitDisplayMode (GLUT_SINGLE | GLUT_RGB);
	glutInitWindowSize (800, 800);
	glutInitWindowPosition (500, 100);
	glutCreateWindow ("Circle");
	init();
    glutDisplayFunc(display);
	glutMainLoop();
	return 0;
}


\end{lstlisting}

\newpage
%-------------------------------------
\section{Output}
\begin{figure}[!h]
    \centering
    
\includegraphics[width=0.7\textwidth]{output.png}
    \caption{Output of Flag Using OpenGL}
\end{figure}

%-------------------------------------
\section{Discussion}
This experiment illustrates how basic OpenGL drawing primitives can be combined to create a meaningful graphical object. The use of an orthographic projection simplifies positioning and scaling in a 2D space. Implementing a custom circle-drawing function enhances understanding of how curves are approximated using line segments in computer graphics. \\

The experiment also demonstrates effective use of color composition and coordinate transformation. Overall, it provides practical insight into structured graphical programming and reinforces theoretical concepts taught in computer graphics lectures.


%-------------------------------------
\section{Conclusion}
In this lab experiment, a 2D flag was successfully drawn using OpenGL and GLUT. The task improved understanding of OpenGL initialization, coordinate systems, geometric primitives, and custom shape generation. This experiment serves as a strong foundation for advanced graphics topics such as transformations, animations, and interactive graphics applications.





\end{document}